% Chapter 5: Context as Currency
\chapter{Context as Currency}
\label{ch:context}

\begin{itshape}
Three hours into a session, Claude starts repeating itself.
It forgets a rule you stated 20 minutes ago.
It suggests an approach you already rejected.
The code quality has noticeably dropped since the session started.
\end{itshape}

In \cref{ch:first-session}, you learned two context commands: \code{/clear}
and \code{/compact}. This chapter explains \emph{why} context degrades,
gives you a full toolkit for managing it, and introduces patterns
for cross-session persistence.

\section{Why Context Degrades}
\label{sec:context-degradation}

\marginnote[Official]{Performance degrades as context fills.
Use \code{/clear} between unrelated tasks. Use \code{/compact}
for selective summarization.\sourceurl{https://code.claude.com/docs/en/best-practices}}

\begin{whybox}[The Non-Linear Cost of Context]
Context degradation is not gradual---there is a noticeable threshold
where Claude begins to lose track of earlier instructions.
Every token of context spent on irrelevant information
(old errors, abandoned approaches, unrelated tasks)
is a token unavailable for the task at hand.

The mechanism: as the context window fills, Claude's attention
distributes across more tokens. Critical instructions
(your CLAUDE.md rules, the current task description) compete
for attention with accumulated noise. Past a threshold,
signal-to-noise drops sharply.
\end{whybox}

% -------------------------------------------------------------------
\section{The Full Context Toolkit}
\label{sec:context-toolkit}

\begin{description}
  \item[\code{/clear}] Wipe context and start fresh. Use between unrelated tasks.
    Cost: 30 seconds. Benefit: full attention on the next task.

  \item[\code{/compact}] Summarize conversation while preserving key decisions.
    Less aggressive than \code{/clear}; retains important context.
    Use when a long task genuinely needs accumulated context but the
    window is getting large.

  \item[\code{/rewind}] Checkpoint and roll back to a previous state.
    Every action Claude takes creates an automatic checkpoint.
    Double-tap \code{Esc} or run \code{/rewind} to open the rewind menu:
    restore conversation only, restore code only, restore both,
    or summarize from a selected message.
    Checkpoints persist across sessions---you can close your terminal
    and still rewind later.

  \item[\code{/context}] Visualize current context usage.
    Shows how much of the window is consumed and by what.

  \item[Subagents] Delegate context-heavy research to isolated agents.
    Results return to your session; the research context does not.
    See \cref{ch:agents} for the full treatment.
\end{description}

\margintip{Watch the token counter. When it approaches 60--70\% of the window,
proactively \code{/compact} or \code{/clear}.}

% -------------------------------------------------------------------
\section{The Two-Failure Rule (Revisited)}
\label{sec:two-failure}

This rule was introduced in \cref{ch:first-session}.
Now you understand \emph{why} it works:

\practitioner{After two failed corrections, the context contains:
the original error ($\sim$500 tokens), your first correction ($\sim$200 tokens),
Claude's apology and retry ($\sim$800 tokens), your second correction ($\sim$200 tokens),
and Claude's second attempt ($\sim$800 tokens).
That is $\sim$2,500 tokens of noise---nearly all of it encoding
failure patterns rather than the actual task.}

\code{/clear} discards 2,500 tokens of noise.
A better initial prompt (informed by what went wrong)
takes $\sim$200 tokens and produces a better result.

\subsection{Customizing Compaction}

\official{You can control what survives compaction by adding instructions
to CLAUDE.md:}

\begin{minted}{markdown}
## Compaction Rules
When compacting, always preserve:
- The full list of modified files
- Any test commands and their results
- Architecture decisions from this session
\end{minted}

For partial compaction, use \code{Esc + Esc} or \code{/rewind},
select a message checkpoint, and choose \textbf{Summarize from here}.
This condenses messages from that point forward while keeping
earlier context intact---more surgical than \code{/compact}.

% -------------------------------------------------------------------
\section{Session Management}
\label{sec:sessions}

\marginnote[Official]{Named sessions: \code{/rename}, \code{--continue},
\code{--resume}, \code{--from-pr}.\sourceurl{https://code.claude.com/docs/en/overview}}

\official{Sessions persist across terminal restarts.} Key commands:

\begin{description}
  \item[\code{claude --continue}] Resume the most recent session.
  \item[\code{claude --resume}] Pick from a list of previous sessions.
  \item[\code{claude --from-pr 123}] Start a session with PR context pre-loaded.
  \item[\code{/rename "feature-auth"}] Give descriptive names for later retrieval.
\end{description}

\subsection{The CURRENT\_WORK.md Pattern}

\practitioner{Maintain a \code{CURRENT\_WORK.md} file at project root.
It provides 30-second context resume after any interruption.}

\begin{fullwidth}
\begin{beforeafter}
\before
\begin{minted}{text}
[No file. Developer spends 10 minutes re-discovering context.]
\end{minted}

\after
\begin{minted}{markdown}
## Right Now
Refactoring feature engineering to use
sklearn Pipelines.

## Why
Current feature code leaks test data during
cross-validation.

## Next Step
Write tests verifying no leakage across CV folds.

## Context When I Return
- Pipeline code in src/features/pipeline.py
- 3 of 7 transformers converted, 4 remaining
- Tests in tests/features/ — existing tests pass
\end{minted}
\end{beforeafter}
\end{fullwidth}

\margintip{Update CURRENT\_WORK.md at the start and end of every session.
60 seconds invested saves 10 minutes of re-discovery.}

% -------------------------------------------------------------------
\section{Auto-Memory and Cross-Session Persistence}
\label{sec:memory}

\marginnote[Official]{Claude records and recalls memories across sessions
automatically. Memory files persist in \code{\textasciitilde/.claude/projects/*/memory/}.\sourceurl{https://code.claude.com/docs/en/memory}}

\official{Claude maintains auto-memory that persists across conversations.}
It records patterns, decisions, and lessons learned, making them available
in future sessions without explicit re-prompting.

\subsection{What to Store in Memory}

\begin{description}
  \item[Stable patterns] Confirmed conventions that apply across sessions
  \item[Architecture decisions] Why you chose X over Y (not just what)
  \item[Hard-won lessons] Bugs, workarounds, gotchas that took time to discover
  \item[Known issues] Pre-existing problems to avoid re-investigating
\end{description}

\subsection{What NOT to Store}

\begin{description}
  \item[Transient state] Current task details, in-progress work
  \item[Unverified conclusions] Patterns seen once, not confirmed
  \item[Duplicates of CLAUDE.md] Memory supplements configuration, does not replace it
\end{description}

\practitioner{The test for whether something belongs in memory:
``Will this still be true and useful in three months?''
If yes, store it. If uncertain, wait for confirmation.}

% -------------------------------------------------------------------
\section{Plan Documents for Large Tasks}
\label{sec:plans}

\marginnote[Official]{Use Plan Mode (Shift+Tab) for read-only analysis
before implementation.\sourceurl{https://code.claude.com/docs/en/overview}}

\official{The recommended workflow for complex tasks:
Explore $\rightarrow$ Plan $\rightarrow$ Implement $\rightarrow$ Commit.}

Plan Mode lets Claude analyze your codebase without modifying anything.
Use it for architecture discussions, impact analysis, and strategy decisions.

\practitioner{For tasks exceeding one hour or 500 lines of changes,
create a plan document:}

\begin{enumerate}
  \item \textbf{Active} --- currently being planned, open questions remain
  \item \textbf{Ready} --- plan approved, implementation can begin
  \item \textbf{Implemented} --- all phases complete
  \item \textbf{Archived} --- moved to \code{docs/plans/archived/}
\end{enumerate}

Every plan document should include a ``Decisions Made'' section
documenting: why specific implementations were chosen,
alternatives considered but rejected, assumptions made,
and tradeoffs accepted. This section is invaluable for postmortem analysis
when a decision turns out to be wrong.

\marginnote[Cross-Ref]{For collaborative plan drafting and ADRs,
see \cref{ch:thinking-partner}.}

% -------------------------------------------------------------------
\section{Quick Reference}

\begin{itemize}
  \item Context degrades non-linearly---manage it actively
  \item \code{/clear} between tasks, \code{/compact} within long tasks
  \item \code{/context} to visualize usage, \code{/rewind} to undo missteps
  \item Two-failure rule: \code{/clear} + better prompt beats a third correction
  \item CURRENT\_WORK.md for cross-session context at project root
  \item Auto-memory for cross-session patterns
  \item Plan Mode (Shift+Tab) before complex implementations
\end{itemize}

% -------------------------------------------------------------------
\begin{trybox}
\begin{enumerate}
  \item Create a \code{CURRENT\_WORK.md} for your active project
    with Right Now, Why, Next Step, and Context sections.
    Include which pipeline step or notebook you are working on.
  \item In your next session, check \code{/context} after 30 minutes.
    How much of the window is consumed?
  \item Track correction rounds for one day. If any task exceeds two,
    note what a better initial prompt would have been.
\end{enumerate}
\end{trybox}
